%!TEX root = tfg.tex
\chapter{Objetivos}

\begin{quotation}[Writer]{Lewis Carroll (1832--1898)}
If you don’t know where you are going, any road will get you there.
\end{quotation}

\begin{abstract}
Este capítulo detalla los motivos que impulsan el desarrollo de Via Inferni, 
analizando la oportunidad de mercado en el género roguelike y el público al que va 
dirigido. Asimismo, se definen de forma precisa los objetivos técnicos y funcionales 
que guiarán la construcción del videojuego bajo una metodología de Ingeniería del Software, 
asegurando la creación de un sistema modular, escalable y rejugable.
\end{abstract}

\section{Motivación}
El desarrollo de este proyecto nace de la intersección entre la literatura clásica y 
las tendencias actuales de la industria del videojuego independiente. El género 
\textit{roguelike} ha demostrado una gran viabilidad comercial debido a su alta 
rejugabilidad y costes de producción contenidos. Sin embargo, existe una oportunidad 
para integrar narrativas profundas, como la 
\textit{Divina Comedia}, en mecánicas de juego procedurales.

Más allá del análisis de mercado, este trabajo surge del interés por experimentar 
el ciclo de vida completo del desarrollo de un software de entretenimiento, desde la 
concepción de la idea hasta su implementación y validación. El desafío de transformar 
conceptos abstractos y reglas de diseño en un producto interactivo funcional supone un 
reto intelectual y técnico que motiva la realización de este proyecto. Como estudiantes 
de Ingeniería Informática a los que nos encanta el sector de los videojuegos, 
este TFG representa la oportunidad de aplicar todos los conocimientos 
adquiridos durante el grado en un contexto real, enfrentando problemas complejos de 
arquitectura de software, optimización y gestión de proyectos.

La hipótesis central es que la estructura de los nueve círculos del infierno es ideal 
para un diseño de niveles ascendentes en dificultad. Cada círculo funciona como un nivel 
temático con ambientación, enemigos y desafíos propios, culminando en  el enfrentamiento contra 
el jefe final: Satán. El público objetivo son jugadores aficionados 
a retos de habilidad y la estética \textit{pixel art}, que buscan experiencias de juego 
cortas pero intensas.

\section{Objetivos}

\subsection{Objetivo general}

El objetivo principal de este Trabajo de Fin de Grado es el desarrollo de un 
prototipo funcional del videojuego \textit{Via Inferni}, inspirado en la 
\textit{Divina Comedia}, aplicando de forma sistemática principios de 
Ingeniería del Software. En particular, se busca integrar prácticas de gestión 
de proyectos basadas en PMBOK con un enfoque de desarrollo iterativo apoyado en 
Feature-Driven Development (FDD), con el fin de evaluar su idoneidad en el 
contexto del desarrollo de videojuegos.

\subsection{Objetivos específicos}

Para alcanzar este objetivo general, se establecen los siguientes objetivos 
específicos, organizados en tres dimensiones: gestión del proyecto, desarrollo 
técnico y validación de resultados.

\subsubsection{Objetivos de gestión del proyecto}

\begin{description}
    \item \textbf{OG1. Planificación del proyecto.} Definir el alcance, cronograma 
    y recursos del proyecto siguiendo las áreas de conocimiento de PMBOK.

    \item \textbf{OG2. Gestión iterativa del desarrollo.} Adaptar la metodología 
    Feature-Driven Development (FDD) al contexto del desarrollo de videojuegos, 
    estructurando el trabajo en iteraciones basadas en funcionalidades 
    incrementales.

    \item \textbf{OG3. Seguimiento y control.} Establecer mecanismos de 
    monitorización del progreso mediante métricas de cumplimiento de hitos, 
    esfuerzo y evolución del producto.

    \item \textbf{OG4. Gestión de riesgos y calidad.} Identificar riesgos 
    potenciales del proyecto y definir estrategias de mitigación, así como 
    criterios de calidad para el producto final.
\end{description}

\subsubsection{Objetivos técnicos del producto}

\begin{description}
    \item \textbf{OT1. Diseño de una arquitectura desacoplada.} Implementar una 
    estructura software basada en el patrón \textit{Observer} y el uso de 
    \textit{ScriptableObjects} para minimizar las dependencias entre sistemas.

    \item \textbf{OT2. Generación procedural de niveles dinámicos.} Desarrollar un
    sistema de generación basado en el algoritmo \textit{Drunkard’s Walk} 
    modificado, garantizando la conectividad y diversidad de las salas.

    \item \textbf{OT3. Implementación del sistema de dualidad de personajes.} 
    Programar la alternancia entre Dante y Virgilio mediante un objeto de 
    persistencia que permita mantener el estado de la partida.

    \item \textbf{OT4. Desarrollo de un ecosistema de enemigos.} Implementar un 
    sistema de comportamiento basado en máquinas de estados finitos para gestionar 
    múltiples tipos de enemigos y jefes.

    \item \textbf{OT5. Sistema de progresión y economía.} Diseñar e implementar un 
    sistema de inventario, objetos y tienda que contribuya a la rejugabilidad del 
    videojuego.

    \item \textbf{OT6. Interfaz y experiencia de usuario.} Desarrollar una interfaz 
    funcional que incluya elementos como minimapa dinámico y gestión de información 
    relevante para el jugador.
\end{description}

\subsubsection{Objetivos de validación}

\begin{description}
    \item \textbf{OV1. Evaluación del proceso de desarrollo.} Analizar la adecuación 
    del uso combinado de PMBOK y FDD en el contexto del proyecto realizado.

    \item \textbf{OV2. Aseguramiento de la calidad.} Diseñar y ejecutar un plan de 
    pruebas que permita validar la estabilidad del sistema y detectar errores.

    \item \textbf{OV3. Balanceo y jugabilidad.} Realizar sesiones de 
    \textit{playtesting} para ajustar la curva de dificultad y mejorar la 
    experiencia del usuario.

    \item \textbf{OV4. Análisis de resultados.} Evaluar los resultados obtenidos en 
    términos de cumplimiento de objetivos, calidad del producto y viabilidad del 
    enfoque metodológico adoptado.
\end{description}